% Document class `report-template` accepts either project-plan or final-report option in
% []. This will change the title page as necessary.
% \documentclass[project-plan]{report-template}
\documentclass[final-report]{report-template}

% Packages I use in my report.
\usepackage{graphicx}
\usepackage{amsmath}
\usepackage{mathtools}
\usepackage{amsfonts}
\usepackage{blindtext}
\usepackage{appendix}

% reusable maths componets
\newcommand{\utt}{\frac {\partial^2 u} {\partial t^2}}
\newcommand{\lap}{\frac {\partial^2 u} {\partial x^2} + \frac{\partial^2 u}{\partial y^2}}
\newcommand{\wave}{\frac{\partial^2 u}{\partial t^2} - c^2\nabla^2 u = s}

% Directory where I saved my figures.
\graphicspath{{./figures/}}

% Metadata used for the title page.
\university{Imperial College London}
\department{Department of Earth Science and Engineering}
\course{MSc in Applied Computational Science and Engineering}
\title{Solving the Wave Equation by Optimising a Discrete Loss}
\author{Luke Dinsdale}
\email{luke.dinsdale22@imperial.ac.uk}
\githubusername{ada-ld2022}
\supervisors{Dr Lluis Guasch Batalla\\
             Dr Ben Moseley\\
             Dr Carlos Cueto Mondejar}
\repository{https://github.com/ese-ada-lovelace-2025/irp-ld2022}

\begin{document}

\maketitlepage  % generate title page

% Abstract
\section*{Abstract}
\label{sect:abstract}

The numerical modelling of waves has broad applications across science and engineering. There are many ways to represent waves computationally, but each requires discretisation, relying on discretisation schemes like the finite difference method. Discretisation introduces errors that can compound over time, particularly when applied to the wave equation. Recently, the Optimising a Discrete Loss (ODIL) framework has been proposed as an efficient method for producing accurate approximations of the solution of Partial Differential Equations (PDE) by means of minimising a discrete loss function. ODIL does not require a time stepping approach, unlike classical numerical methods, which may offer accuracy and efficiency benefits when solving the wave equation. This work aims to explore the accuracy and performance characteristics of the ODIL framework when applied to the wave equation.

% Introduction section
\section{Introduction}
\label{sect:intro}

\section{Background}
\label{sect:backg}

\subsection{The Acoustic Wave Equation}
\label{subsect:wave_eq}

Mechanical wave propagation through a medium can be modelled using the linear acoustic wave equation. Assuming a lossless, isotropic medium, the pressure field $u = u(t,x,y)$ satisfies

\begin{equation}
    \wave,
    \label{eq:wave}
\end{equation}

where $c = c(x,y)$ is the spatially varying wave speed, $\nabla^2$ denotes the two-dimensional Laplacian operator,

\begin{equation}
    \nabla^2
    =
    \frac{\partial^2}{\partial x^2}
    +
    \frac{\partial^2}{\partial y^2},
    \label{eq:lap}
\end{equation}

with time $t$ and the spatial coordinates $(x,y)$. $s = s(t)$ is the acoustic source modelled as a time shifted Ricker wavelet \cite{rickerWaveletFunctionsTheir1944}, 

\begin{align}
    s(t) &= \left( 1 - 2\alpha \right)e^{ -\alpha}, \label{eq:ricker}\\
    \alpha &= \pi^2 f^2_0(t - t_0)^2,
\end{align}

where $f_0$ is the peak frequency and $t_0$ is the time shift that centres the pulse.

\subsubsection{Initial and Boundary Conditions}
\label{subsect:conditions}

The pressure field is assumed to be initially at rest, corresponding to zero initial pressure and zero initial particle velocity prior to excitation by the acoustic source. Correspondingly, we subject the wave equation to the initial conditions

\begin{align}
    u(0, x, y) &= 0 \label{eq:u0} \\
    \frac{\partial}{\partial t}u(0, x, y) & = 0 \label{eq:ut0}
\end{align}

Although acoustic propagation occurs in an effectively unbounded medium, numerical simulations must be performed on a finite computational domain. Appropriate absorbing boundary conditions are therefore required to minimise non-physical reflections from the artificial boundaries. We impose second order Higdon conditions \cite{higdonNumericalAbsorbingBoundary1987a}, which provide improved absorption of outgoing waves over a wider range of incidence angles than first order absorbing conditions while remaining significantly less computationally expensive than perfectly matched layers. They therefore provide a suitable balance between accuracy and computational cost for the present work.

In two dimensions, the second order Higdon absorbing boundary conditions read

\begin{equation}
    \left( \frac{\partial}{\partial t} + c_1 \frac{\partial}{\partial n} \right)
    \left( \frac{\partial}{\partial t} + c_2 \frac{\partial}{\partial n} \right) u = 0,
    \label{eq:higdon}
\end{equation}

where the normal derivative

\begin{equation}
    \frac {\partial u} {\partial n} = \nabla u \cdot \mathbf{n},
    \label{eq:dudn}
\end{equation}

is the spatial derivative of the pressure field in the direction perpendicular to the boundary with $\mathbf n$ as the outward unit normal. The parameters $c_1$ and $c_2$ represent the propagation velocities of outgoing waves that the boundary condition is designed to absorb. For a homogeneous medium, these are commonly chosen to be

\begin{equation}
    c_n = c \cos(\theta_n),
    \label{eq:phasevel}
\end{equation}

where $\theta_n$ is the angle of incidence relative to the boundary normal. 

In Appendix \ref{app:hig_exp}, we show that Equation \eqref{eq:higdon} can be written as

\begin{equation}
    \frac{\partial^2 u}{\partial t^2} + 
    (c_1 + c_2) \frac{\partial^2 u}{\partial t \partial n } + 
    c_1 c_2 \frac{\partial^2 u}{\partial n^2} = 0.
    \label{eq:higdon_expanded}
\end{equation}

where mixed partial derivatives have been combined using the commutativity of partial differentiation. We use this form because it is much easier to discretise, since each derivative can be evaluated independently according to predefined finite difference operators which we discuss in \ref{subsubsect:disc_op}.

\subsection{Optimising a Discrete Loss}

Unlike conventional finite difference methods, which solve the wave equation by explicitly marching forward in time, ODIL reformulates the governing PDE as an optimisation problem over the complete discrete space time domain. The discrete field itself is treated as the optimisation variable, and the solution is obtained by minimising a loss function constructed from the residual of the discretised governing equations.

In Section \ref{subsect:discretisation}, we show that the wave equation can be written as a discrete linear system

\begin{equation}
    A \bf u = s,
\end{equation}

where $A$ encodes the action of discrete approximations of the derivative operators and the aforementioned initial and boundary conditions, $\bf u$ (the optimisation variable) is a discrete representation of our pressure field, and $\bf s$ is our modelled source. Accordingly, the discrete residual can be written

\begin{equation}
    r(\mathbf{u}) = A \bf u - s.
    \label{eq:residual}
\end{equation}

In the ODIL framework, the optimisation variable comprises the entire discrete space-time pressure field rather than the solution at a single time level, as is the case in conventional time stepping methods. Hence, we choose to write the wave equation in the form of an all-at-once system as above.

Under this formulation, the discrete loss function $\mathcal{L}$ is 

\begin{equation}
    \mathcal{L}(\mathbf{u}) = \frac 1 2 \|\mathbf{r}(\mathbf{u})\|^2,
    \label{eq:cont_loss}
\end{equation}

so the residual vanishes when the discrete wave equation is satisfied. The forward problem we seek to solve is therefore the unconstrained optimisation problem

\begin{equation}
    \min_{\mathbf{u}}\mathcal{L}(\bf u).
\end{equation}

Then it follows that the minimiser is a discrete pressure field $\bf u$ for which the residual is minimised and, in the ideal case, vanishes entirely, implying satisfaction of the discretised wave equation.

\section{Methodology}
\label{sect:methods}

\subsection{Finite Difference Discretisation}
\label{subsect:discretisation}

To solve the wave equation numerically, the continuous domain is discretised into a finite set of spatial and temporal grid points. The resulting discrete approximation replaces the continuous differential operators with finite difference operators, allowing the wave equation to be expressed as a finite dimensional system suitable for numerical solution.

The computational domain is discretised using a uniform Cartesian grid with $N_t,\ N_x,\ N_y$ grid nodes in the temporal and spatial directions, respectively. The spatial coordinates are then defined as

\begin{equation}
    x_i \coloneqq i\Delta x, \quad y_j \coloneqq j\Delta y
    \label{eq:spatdisc}
\end{equation}

and temporal coordinates

\begin{equation}
    t_n \coloneqq n\Delta t,
    \label{eq:tempdisc}
\end{equation}

with $n = 0, 1, \dots N_t - 1$ (and analogously for $x$ and $y$) and $\Delta x$, $\Delta y$, and $\Delta t$ denote the spatial and temporal grid spacings, respectively.

Under this discretisation, our discrete pressure field $u^n_{i,\, j}$ becomes an approximation of the continuous pressure field $u$:

\begin{equation}
    u^n_{i,\, j} \approx u(t_n, x_i, y_j),
    \label{eq:approxu}
\end{equation}

The spatial and temporal second derivatives in Equation \eqref{eq:wave} are discretised with standard centred finite differences. In practice we employ higher-ordered schemes (4th, 6th, 8th) detailed in Appendix \ref{app:central_diff}, but for the purpose of this discussion we will make reference to the 2nd order central differences,

\begin{align}
    \utt &\approx \frac {u^{n+1}_{i,\, j} - 2u^{n}_{i,\, j} + u^{n-1}_{i,\, j}} {\Delta t^2} \label{eq:uttcd}\\[2pt]
    \frac{\partial^2 u}{\partial x^2} &\approx \frac {u^{n}_{i + 1,\, j} - 2u^{n}_{i,\, j} + u^{n}_{i-1,\, j}} {\Delta x^2} \label{eq:uxxcd}\\[2pt]
    \frac{\partial^2 u}{\partial y^2} &\approx \frac {u^{n}_{i,\, j+1} - 2u^{n}_{i,\, j} + u^{n}_{i,\, j-1}} {\Delta y^2} \label{eq:uyycd}
\end{align}

\subsubsection{Matrix Representation of the Discrete Operators}
\label{subsubsect:disc_op}

Unlike conventional time stepping approaches, which solve for the wavefield sequentially in time, ODIL formulates the discretised PDE as a global space time problem. The unknown wavefield is represented as a single vector containing all spatial and temporal degrees of freedom. Consequently, the differential operators are naturally represented as sparse matrices acting on this global vector.

The finite difference discretisation transforms the continuous wavefield into a set of unknown values defined on the discrete space-time grid. These values are collected into a vector representation,

\begin{equation}
    \mathbf{u} \coloneqq \left[\left(u^0\right)^T,\ \left(u^1\right)^T,\ \dots,\left(u^{N_t - 1}\right)^T\right]^T.
    \label{eq:u_vec}
\end{equation} 

Here the spatial indices are flattened into a single spatial degree of freedom index where $u^n \in \mathbb{R}^{N_x N_y}$ is the pressure field across space at the $n$-th time level.

We define the temporal second derivative operator $D_{tt} \in \mathbb{R}^{N_t \times N_t}$ and the second spatial derivative operators $D_{xx} \in \mathbb{R}^{N_x \times N_x}$ and $D_{yy} \in \mathbb{R}^{N_y \times N_y}$. These matrices represent the action of the finite difference approximations of second derivatives on $\mathbf{u}$. For illustration, the second order centred approximation to the temporal second derivative gives the sparse matrix

\begin{equation}
    D_{tt} = \frac{1}{\Delta t^2}\begin{bmatrix}
        -2 & 1 & 0 & \cdots & 0 \\
        1 & -2 & 1  & \ddots & \vdots\\
        0 & 1 & -2 & \ddots  & 0 \\
        \vdots & \ddots & \ddots & \ddots & 1\\
        0 & \cdots & 0 & 1 & -2
    \end{bmatrix}.
    \label{eq:DTT_mat}
\end{equation}

The finite difference stencil therefore appears as the non-zero diagonals of the operator matrix. Higher order schemes retain the same sparse structure but introduce additional off diagonal entries corresponding to the wider stencil. Note here that, in the absence of boundary modifications, the finite difference derivative matrices possess a Toeplitz structure, illustrated in Figure \ref{fig:sparsity}.

The discrete Laplacian operator is constructed using Kronecker products as

\begin{equation}
    L_{xy}
    \coloneqq
    D_{xx}\otimes I_y
    +
    I_x\otimes D_{yy},
    \label{eq:lap}
\end{equation}

where $L_{xy}\in\mathbb{R}^{N_xN_y\times N_xN_y}$ acts on the vectorised spatial wave field. Here, $I_x \in \mathbb{R}^{Nx\times Nx}$ and $I_y \in \mathbb{R}^{Ny\times Ny}$ are identity matrices.   

To write the wave equation as a linear system, we apply Kronecker products to form a global matrix operator $A \in \mathbb{R}^{N_t N_x N_y \times N_t N_x N_y}$ that represents the action of the second temporal derivative and Laplacian across the full space-time vector,

\begin{equation}
    A \coloneqq D_{tt} \otimes I_{xy} - I_t \otimes \left(C^2 L_{xy}\right),
    \label{eq:A_op}
\end{equation}

where $I_t \in \mathbb{R}^{Nt \times N_t}$ and $I_{xy} \in \mathbb{R}^{N_x \times N_y}$ are the identity matrices and $C \in \mathbb{R}^{N_x N_y \times N_x N_y}$ is a diagonal matrix containing the spatially varying wave speed values on the computational grid:

\begin{equation}
    C = \text{diag}(c_1, c_2, \dots, c_{N_xN_y}).
\end{equation}

The Kronecker product lifts the local finite difference stencil into a global sparse space time operator, where each temporal coupling is represented by a block matrix acting on the spatial degrees of freedom, for instance,

\begin{equation}
    D_{tt} \otimes I_{xy} = \begin{bmatrix}
        -2I_{xy} & I_{xy} & 0 & \cdots \\
        I_{xy} & -2I_{xy} & I_{xy}  & \cdots\\
        0 & I_{xy} & -2I_{xy} & \cdots \\
        \vdots & &  & \ddots \\
    \end{bmatrix}.
    \label{eq:dtt_kron}
\end{equation}

The resulting space time operator has a Block Toeplitz with Toeplitz Blocks (BTTB) structure, where each block corresponds to a time level, diagonal blocks represent the current time level, and off-diagonal blocks couple neighbouring time levels. The BTTB structure is highlighted in Figure \ref{fig:sparsity}.

Under this discretisation, the full system can be expressed as

\begin{equation}
    A\mathbf{u} = \mathbf{s}.
    \label{eq:Aus}
\end{equation}

This equation represents the discrete PDE residual evaluated over the complete space time domain. Unlike conventional time stepping methods, in the ODIL framework the wavefield $\bf u$ is optimised simultaneously across all temporal and spatial degrees of freedom, which is why we choose to form the wave equation into an all-at-once system. Note that, in practice, the action of the global operator is evaluated matrix free through a sequence of sparse matrix-vector products corresponding to the temporal and spatial operators. This avoids storing the full space time operator, whose size scales with the total number of space-time degrees of freedom. This sparsity is illustrated in Figure \ref{fig:sparsity}.

\begin{figure}[!h]
    \centering
    \includegraphics[width=\linewidth]{latex/figures/fig1_op_structure.png}
    \caption{Sparsity plots of the operator matrices $D_{xx},\ D_{yy}$, and $D_{tt}$ (a - c) and the global operator matrix $A$ (d) in the discretised wave equation for a second order centred difference. Blue squares represent the non-zero entries in the matrices, and red boxes indicate the BTTB structure of the global operator. The inset illustrates the block structure: off-diagonal blocks couple neighbouring time levels and the central block represents the current time with the spatial operator terms folded in. The huge number of zero entries in the global operator matrix motivates a matrix free approach and Compressed Sparse Row format for the stored constituents.}
    \label{fig:sparsity}
\end{figure}

With the wave equation written as a linear system, we arrive at the discrete ODIL residual introduced in Equation \eqref{eq:residual},

\begin{equation}
    r(\mathbf{u}) = A\mathbf{u} - \mathbf{s},
    \label{eq:Aus}
\end{equation}

which we seek to minimise.
\newpage
\subsubsection{Discretisation of the Initial and Boundary Conditions}
\label{subsubsect:disc_conditions}

The formulation above describes the interior discretisation of the wave equation. Initial and boundary conditions are incorporated by modifying the corresponding rows of the discrete operator, such that the complete system represents both the PDE residual and constraint equations within the same global formulation.

The initial pressure field is imposed directly by setting the solution at the initial time level to zero. The zero initial velocity condition is enforced using a finite difference approximation of the first temporal derivative at $t=0$,

\begin{align}
    u(0,\ x_i,\ y_j) = u^0 & = 0 \label{eq:disc_uo} \\[2pt]
    \frac {\partial} {\partial t} u(0,\ x_i,\ y_j) \approx \frac{u^1 - u^0}{\Delta t} &= 0 \label{eq:disc_ut0}
\end{align}

This approximation simply reduces to rearranging Equation \eqref{eq:disc_ut0} such that it reads

\begin{equation}
    u^1 = u^0 = 0 \quad \Longrightarrow \quad u^1 = 0,
\end{equation}

or, in plain terms, to pinning $u(t_1, x_i, y_j) = 0$.

At the computational boundaries, the interior finite difference stencil is replaced by a discrete approximation of the second-order Higdon absorbing boundary condition (Equation \eqref{eq:higdon_expanded}):

\begin{equation}
    \delta_{tt} u + (c_1 + c_2)\delta_{tn} u + c_1c_2\delta_{nn} u = 0,
    \label{eq:disc_higdon}
\end{equation}

where $\delta_{tt}$ denotes the discrete second derivative operator in time, $\delta_{tn}$ the discrete mixed derivative operator, and $\delta_{nn}$ the discrete second derivative of the normal coordinate. These differentiation operators are approximated by first and second order centred finite differences as in Equations \eqref{eq:uttcd} - \eqref{eq:uyycd}.

\subsubsection{Enforcing the Boundary Condition}
\label{subsubsect:ghost_nodes}

Centred finite difference stencils of order $2g$ require $g$ nodes on either side of the evaluation points and therefore are not directly applicable within $g$ nodes of the domain boundary. Rather than switching to one sided stencils near the boundary, the domain is extended by $g$ ghost nodes on each side, whose values are defined by polynomial extrapolation of the interior solution. Each ghost value is obtained by fitting a Lagrange polynomial through $k = \min(4,\ 2g)$ interior nodes nearest the boundary and evaluating it at the ghost location $x_g$. Specifically, the ghost node value $u_g$ is approximated by

\begin{equation}
    u_g = \sum^{k-1}_{i=0} \ell_i(x_g)u(x_i),
    \label{eq:lagrange}
\end{equation}

where $x_i$ denote the neighbouring interior nodes and

\begin{equation}
    \ell_i(x) = \prod^{k-1}_{j=0 \\ j\neq i} \frac {x - x_j} {x_i - x_j},
    \label{eq:lag_base}
\end{equation}

are the Lagrange basis polynomials.

This approximation is exact for polynomials of degree $k-1$. Consequently, linear extrapolation is employed for second order stencils, while cubic extrapolation for stencils of fourth order and higher. The polynomial degree is capped at 3 to avoid the ill-conditioning associated with high degree extrapolation, which can amplify perturbations in the interior solution through rapidly growing extrapolation weights. 

With cubic extrapolation, the ghost nodes carry an $O(h^4)$ truncation error, matching the boundary accuracy of the fourth order stencils. As a result, the formal accuracy of the fourth order discretisation is preserved. For the sixth and eighth order schemes, the cubic cap formally reduces the order of accuracy at the boundary. However, this trade off was found to provide an improvement in numerical robustness by avoiding the instability associated with high degree extrapolation without a significant decrease in accuracy at the boundary. This approach allows the same high order finite difference stencil to be retained at boundary nodes without introducing additional degrees of freedom into the system.

\subsubsection{Matrix Form of the Boundary Operator}
\label{subsubect:matrix_bc}

Since the extrapolation weights depend only on the grid geometry, the ghost node reconstruction is a linear operation. It can therefore be represented as a sparse matrix

\begin{equation}
    \mathbf{u}_{\text{full}} = G\mathbf{u},
    \label{eq:ghost}
\end{equation}

where $G\in\mathbb{R}^{(n+2g)\times n}$ maps the $n$ interior unknowns to the extended vector containing both interior and ghost nodes. The matrix $G$ consists of an identity block corresponding to the interior nodes and additional rows containing the Lagrange extrapolation weights for the ghost nodes.

Accordingly, in the $x$ direction, we can write

\begin{equation}
    D_x^{(g)} = D_x G, \qquad D_{xx}^{(g)} = D_{xx} G,
\end{equation}

which yields ordinary square difference operators, valid at every node including the boundary, with the extrapolated ghost contribution folded in. We broadcast the discrete operators $\delta_{tt},\ \delta_{tn},\ \delta_{nn}$ of Equation \eqref{eq:disc_higdon} to space-time by means of another Kronecker product as in Equation \eqref{eq:lap}, for instance

\begin{equation}
    D_{tt} = \delta_{tt} \otimes I_{xy},
\end{equation}

and correspondingly for the spatial operators. Then, restricting to the $n_b$ nodes on edge $b$ with a selection matrix $R_b$, we can compute the discrete normal derivative operators as

\begin{equation}
    D_n = R_b D_x^{(g)}, \qquad D_{nn} = R_b D_{xx}^{(g)},
\end{equation}

with $D_x^{(g)}$ replaced by $D_y^{(g)}$ on the horizontal edges, and $\cdot^{(g)}$ denoting the action of the operator at the boundary nodes. Since $\partial_n$ is the outward normal, $D_n$ is signed according to which edge is being treated. Outward is taken to be negative at the left/bottom boundary and positive at the right/top, so $D_n$ flips sign accordingly.

Equation \eqref{eq:disc_higdon} is then assembled for every boundary node and time step,

\begin{equation}
    D_{tt}\mathbf{u}_b + (c_1+c_2) D_{tn}\mathbf{u}_b + c_1c_2 D_{nn}\mathbf{u}_b  = 0,
    \label{eq:mat_higdon}
\end{equation}

where $\mathbf{u}_b$ are the boundary columns of the solution array. To interface with the aforementioned matrix-free philosophy, we apply the boundary conditions as hard constraints by replacing the corresponding rows of the discrete wave operator with the boundary constraint equations.

\subsubsection{Solution}

Gauss Newton linearises this around the current iterate $u_k$ such that

$$r(u+\delta u) \approx r_k + J_k \delta u,$$

notice for our linear system that the Jacobian for the $k$-th iterate is simply $A$. Our loss is then

$$L(u) = \frac 12 \|r_k - A\delta u\|^2,$$

whose unique minimiser is the solution to the normal equations

$$A^TA\delta u = -A^Tr.$$

However since $A$ is nonsingular we can write

$$\underbrace{(A^T)^{-1}A^T}_I \delta u = -\underbrace{(A^T)^{-1}A^T}_I r,$$

which reduces to

$$A\delta u = -r.$$

Thus the unique minimiser of the normal equations in our linear case is merely the solution to the above linear equation. Considering our loss is quadratic, Gauss Newton converges in exactly one step, and minimising our loss becomes a problem of solving the above linear equation.

\subsubsection{Preconditioning}

However, the conditioning number of that system is typically large, and increases superlinearly with the number of time levels in the discretisation. This warrants a time-aware preconditioner to allow the inner solve to converge quickly.

We address this using an $\omega$-circulant preconditioning matrix \cite{danieliAllatonceSolutionLinear}. Our temporal operator is nearly circulant, it is just missing the periodic diagonal edge values. The core idea is that a circulant matrix operator can be trivially diagonalised by its Fourier modes according to

$$D_{tt}^\omega =   F^{-1}\Lambda F,$$

where $\Lambda$ is the diagonal matrix of Fourier eigenvalues. Assuming the solution is composed of a single Fourier mode such that
$$u = e^{i\theta n}$$

\begin{figure}[!h]
    \centering
    \includegraphics[width=0.9\linewidth]{latex/figures/fig2_sv_distribution.png}
    \caption{Sparsity plot of the global operator matrix $A$, including boundary terms (a), and circulant approximation $M$, the preconditioner (b), along with their singular value distributions (c, d). Small singular values in the unpreconditioned operator lead to severe ill-conditioning and unstable solves of the linear system in the Gauss Newton optimisation process, which the preconditioner addresses.}
    \label{fig:svd}
\end{figure}

% Another section
\section{Methods}

\section{Results}

\section{Discussion}

\section{Conclusion}

\appendix
\section{Appendix}
\subsection{Expansion of the Higdon condition}
\label{app:hig_exp}

Starting from the Higdon boundary condition (Equation \eqref{eq:higdon}), define the interior operator

\begin{equation}
    v \coloneqq \left( \frac{\partial}{\partial t} + c_2 \frac{\partial}{\partial n} \right)u,
\end{equation}

using which the Higdon condition can be written

\begin{equation}
    \left( \frac{\partial}{\partial t} + c_1 \frac{\partial}{\partial n} \right)v = 0.
\end{equation}

Applying the outer operator,

\begin{equation}
    \frac{\partial v}{\partial t} + c_1 \frac{\partial v}{\partial n} = 0,
\end{equation}

and substituting $v$ we arrive at

\begin{equation}
    \frac{\partial}{\partial t} \left( \frac{\partial u}{\partial t} + c_2 \frac{\partial u}{\partial n} \right) +
    c_1 \frac{\partial}{\partial n} \left( \frac{\partial u}{\partial t} + c_2 \frac{\partial u}{\partial n} \right) = 0.
\end{equation}

Differentiating each term gives

\begin{equation}
    \frac{\partial^2 u}{\partial t^2} + 
    c_2 \frac{\partial^2 u}{\partial t \partial n } + 
    c_1\frac{\partial^2 u}{\partial n \partial t} +
    c_1 c_2 \frac{\partial^2 u}{\partial n^2} = 0
\end{equation}

and since mixed derivatives commute we can write

\begin{equation}
     \frac{\partial^2 u}{\partial t \partial n } =  \frac{\partial^2 u}{\partial n \partial t },
\end{equation}

and rearrange to find

\begin{equation}
    \frac{\partial^2 u}{\partial t^2} + 
    (c_1 + c_2) \frac{\partial^2 u}{\partial t \partial n } + 
    c_1 c_2 \frac{\partial^2 u}{\partial n^2} = 0,
\end{equation}

which is the result presented in Equation \eqref{eq:higdon_expanded}.

\subsection{Central Difference Schemes}
\label{app:central_diff}

On the discrete computational domain described in Section \ref{subsect:discretisation}, the temporal and spatial derivatives in the wave equation (Equation \ref{eq:wave}) are approximated using standard central difference schemes of orders 2, 4, 6, and 8. For a $2P + 1$ point stencil in $x$, the general form of the centred difference is

\begin{equation}
    \frac {\partial^2 } {\partial x^2}u(t_n, x_i, y_j) \approx \frac {1} {\Delta x^2} \sum^P_{p=-P}c_p u_{i+m},
\end{equation}

which reads analogously in $y$ and $t$. The coefficients $c_p$ are obtained via the usual Taylor series construction and are chosen to achieve second, fourth, and eighth order accuracy. Those coefficients are tabulated below.

\begin{table}[!htp]
\centering
\begin{tabular}{|c|l|l|}
    \hline
    Order & Offsets & Coefficients \\
    \hline
    2nd & $-1,\ 0 ,\ 1$ & $1,\ -2,\ 1$\\[3pt]
    4th & $-2,\ -1,\ 0 ,\ 1,\ 2$ & $-\frac{1}{12},\ \frac{4}{3},\ -\frac{5}{2},\ \frac{4}{3},\ -\frac{1}{12}$\\[3pt]
    6th & $-3,\ -2,\ -1,\ 0 ,\ 1,\ 2$ & $\frac{1}{90},\ -\frac{3}{20},\ \frac{3}{2},\ -\frac{49}{18},\ \frac{3}{2},\ -\frac{3}{20},\ \frac{1}{90}$\\[3pt]
    8th & $-4,\ -3,\ -2\, -1,\ 0 ,\ 1,\ 2,\ 3,\ 4$ & $-\frac{1}{560},\ \frac{8}{315},\ -\frac{1}{5},\ \frac{8}{5},\ -\frac{205}{72},\ \frac{8}{5},\ -\frac{1}{5},\ \frac{8}{315},\ -\frac{1}{560}$ \\[3pt] \hline
\end{tabular}
\caption{Coefficients for centred finite difference approximations of the second derivative.}
\label{tab:fd_coefficients}
\end{table}

\newpage
% References
\bibliographystyle{plain}
\bibliography{references}  % BibTeX references are saved in references.bib

\end{document}          
